% TOP_PROBLEM: {"id": 3256, "title": "Woodall's Conjecture", "category_id": 3, "category": {"id": 3, "name": "graph_theory", "display_name": "Graph Theory", "description": "Problems involving graphs, networks, and their properties.", "slug": "graph-theory", "order_index": 3, "created_at": "2026-07-31T15:26:25.671Z"}, "status": "open", "problem_number": "OPG-179", "set_id": 12}
% FIRST_POSTED: 2026-10-01
% ENTRY_KIND: statement_only
% UnsolvedMath ID 3256; source catalogue code OPG-179
% !TeX program = lualatex
% Definition and sources only; this file is not an LLM solution attempt.
\documentclass[11pt,a4paper]{article}
\usepackage[margin=25mm]{geometry}
\usepackage{fontspec}
\setmainfont{lmroman10-regular.otf}[BoldFont=lmroman10-bold.otf,
 ItalicFont=lmroman10-italic.otf,BoldItalicFont=lmroman10-bolditalic.otf]
\usepackage{amsmath,amssymb,mathtools}
\usepackage{unicode-math}
\setmathfont{latinmodern-math.otf}
\AtBeginDocument{\let\setminus\smallsetminus}
\usepackage{xurl,hyperref}
\hypersetup{colorlinks=true,urlcolor=blue}
\setcounter{secnumdepth}{0}
\setlength{\parindent}{0pt}
\setlength{\parskip}{5pt}
\setlength{\emergencystretch}{3em}
\begin{document}
\section*{Woodall's Conjecture}\label{3256}
{\small UnsolvedMath ID 3256 · OPG-179 · Graph Theory · OpenGarden}
\subsection{Definitions and mathematical statement}
Let $D=(V,A)$ be a finite loopless digraph; parallel arcs, if present, are
treated as distinct arcs. For $\varnothing\ne X\subsetneq V$, define
$\delta^+(X)=\{xy\in A:x\in X,y\notin X\}$ and define $\delta^-(X)$
by reversing the membership conditions. A directed cut is a nonempty
set $\delta^+(X)$ for which $\delta^-(X)=\varnothing$. Thus every arc
crossing this vertex partition goes in the same direction; an arbitrary
outgoing cut with incoming arcs as well is not a directed cut here.

A dijoin is a set $J\subseteq A$ meeting every directed cut. Suppose
$D$ has at least one directed cut, and let
\[
             k=\min\{|C|:C\text{ is a directed cut of }D\}.
\]
Woodall's conjecture asks for $k$ dijoins $J_1,\ldots,J_k$ with
$J_i\cap J_j=\varnothing$ whenever $i\ne j$. They may share incident
vertices; only the arc sets must be disjoint.

Every directed cut must contain a different representative from each
of the $k$ dijoins, so its size gives an immediate upper bound on the
possible number. The conjecture asserts attainment of that upper bound
by integral arc sets. It does not ask for a fractional packing or for
vertex-disjoint subgraphs. One may equivalently partition all arcs into
$k$ dijoins, since unused arcs can be added to any one of them without
destroying the property of meeting every directed cut.
\subsection{Short English statement}
If the smallest directed cut has k arcs, must there be k pairwise arc-disjoint sets each meeting every directed cut?
\subsection{Sources}
\begin{itemize}
\item[S1] ULAM.AI and the UnsolvedMath Contributors, supplied problems.json, record 3256 (OPG-179), OpenGarden collection. Catalogue identity and source transcription; CC BY 4.0.
\url{https://huggingface.co/datasets/ulamai/UnsolvedMath}.
\item[S2] Open Problem Garden, Woodall's Conjecture. Original problem and discussion; the mathematical formulation below is expanded from the supplied source record.
\url{http://www.openproblemgarden.org/op/woodalls_conjecture}.
\end{itemize}
\end{document}
